% ============================================================
% CAPÍTULO 8 — ANEXOS
% ============================================================
\clearpage
\chapter*{Anexos}
\addcontentsline{toc}{chapter}{Capítulo 8. Anexos}

\section{Anexo A. Parámetros del Sistema}
\begin{table}[H]
\caption{Parámetros principales}\label{tab:parametros}
{\singlespacing
\begin{tabularx}{\textwidth}{@{}lX@{}}
\toprule
Parámetro & Valor \\
\midrule
Fotogramas por clip & 16 \\
Ventana temporal máxima & 7,0~s \\
Hueco máximo entre mediciones & 0,5~s \\
Tolerancia de anclaje & 0,10~s \\
Confianza mínima de ojo y hocico & 0,50 \\
Tamaño del recorte & $224 \times 343$ píxeles \\
Extractor visual & ResNet18 preentrenada en ImageNet, congelada \\
Dimensión visual / geométrica & 512 / 9 \\
Proyección geométrica & $9 \rightarrow 64 \rightarrow 128$ \\
Particiones externas & 4, agrupadas por caballo \\
Épocas máximas / paciencia & 100 / 10 \\
Épocas del modelo de producción & 29 \\
Parámetros entrenables & 132\,355 \\
\bottomrule
\end{tabularx}}
\end{table}

\section{Anexo B. Componentes del Software}
\begin{itemize}
  \item \texttt{equinos/catalogo.py}: catálogo de videos por clase.
  \item \texttt{equinos/reproceso.py}: estimación de pose sobre todos los fotogramas.
  \item \texttt{equinos/secuencia.py}: selección temporal de fotogramas.
  \item \texttt{equinos/alineado.py}: alineación de la cara.
  \item \texttt{equinos/dataset.py}: construcción de los recortes.
  \item \texttt{equinos/caracteristicas.py}: extracción con ResNet18.
  \item \texttt{equinos/clasificacion.py}: modelo, entrenamiento y evaluación.
  \item \texttt{equinos/particion.py}: particiones por caballo.
  \item \texttt{equinos/inferencia.py} y \texttt{app\_inferencia.py}: aplicación de inferencia.
  \item \texttt{scripts/evaluar\_v2.py} y \texttt{scripts/evaluar\_lineas\_de\_base.py}: evaluación del modelo y de las líneas de base.
  \item \texttt{tests/}: pruebas automatizadas, entre ellas la de equivalencia entre entrenamiento e inferencia.
\end{itemize}

\section{Anexo C. Resultados por Partición}
\begin{table}[H]
\caption{Entrenamiento, validación y prueba por partición}\label{tab:folds}
{\singlespacing
\begin{tabular}{@{}rrrrr@{}}
\toprule
Partición & Mejor época & \textit{F}\textsubscript{1} entrenamiento & \textit{F}\textsubscript{1} validación & \textit{F}\textsubscript{1} prueba \\
\midrule
1 & 28 & 0,973 & 0,700 & 0,718 \\
2 & 26 & 0,994 & 0,766 & 0,769 \\
3 & 20 & 0,969 & 0,641 & 0,677 \\
4 & 40 & 0,981 & 0,742 & 0,795 \\
\midrule
Media & 28,5 & 0,979 & 0,712 & 0,740 \\
\bottomrule
\end{tabular}}
\notatabla{\textit{F}\textsubscript{1} macro. El entrenamiento se evalúa sin \textit{dropout} con los pesos de la mejor época.}
\end{table}

\section{Anexo D. Variantes Evaluadas}
\begingroup\singlespacing\small
\setlength{\LTleft}{0pt}
\begin{longtable}{@{}p{0.27\textwidth}p{0.42\textwidth}p{0.23\textwidth}@{}}
\caption{Modelos y variantes evaluados}\label{tab:modelos}\\
\toprule
Modelo & Descripción & Resultado \\
\midrule
\endfirsthead
\toprule
Modelo & Descripción & Resultado \\
\midrule
\endhead
Un fotograma y voto & Regresión logística por fotograma y voto entre 16 predicciones. & 0,660 sobre 153 clips. \\
Modelo de promedio & Proyección compartida y promedio de 16 fotogramas. & 0,704 sobre 153 clips. \\
GRU & Recurrente, sensible al orden. & No superó a los invariantes. \\
Atención con posición & Atención con codificación posicional. & Sin mejora consistente. \\
CNN--LSTM & ResNet18 seguida de LSTM. & No integrada. \\
MIL visual & MIL sólo con características visuales. & 0,686. \\
MIL multimodal & Visual, geometría y atención. & Modelo final: 0,740. \\
Ajuste fino parcial & Última etapa de la ResNet18 entrenable. & 0,747; no adoptado. \\
Canal geométrico ampliado & Más puntos anatómicos. & Peor o más inestable. \\
Puntos de boca y ollares & Puntos faciales adicionales. & Sin mejora. \\
\bottomrule
\end{longtable}
\endgroup
